\chapter*{数列专题练习}
\begin{enumerate}
    \item 求等差数列的通项公式、项和公差
    \begin{enumerate}
        \item 若等差数列$\{a_n\}$的前三项依次为：$a-1, a+1, 2a+3$，则此等差数列的通项$a_n = $\blkda{$2n-3$}
        \begin{jieda}
            $(a-1) + (2a+3) = 2(a+1)$，解得$a = 0$，故$\{a_n\}$的前三项依次为：$-1, 1, 3$，公差为$2$，通项公式为$2n-3$.
        \end{jieda}
        \item 首项为$-24$的等差数列从第10项开始为正数，则公差$d$的取值范围是\blkda{$(\dfrac{8}{3}, 3]$}
        \begin{jieda}
            根据条件可以列出$\left \{ \begin{array}{l}
                 -24 + 8d \leq 0  \\
                 -24 + 9d > 0
            \end{array}\right.$，解得$d \in (\dfrac{8}{3}, 3]$
        \end{jieda}
        \item 若等差数列$18, 15, 12, \cdots$的前$n$项和最大，则$n = $\blkda{6或7}
        \begin{jieda}
            首先可以写出通项公式$a_n = 21-3n$，则易知$a_7 = 0$，$\{a_n\}$为严格减数列，故6或7满足条件.
        \end{jieda}
        \item 等差数列$\{a_n\}$为严格增数列，已知$a_3 + a_6 + a_9 = 12, a_3a_6a_9 = 28$，则此等差数列的通项$a_n = $\blkda{$n-2$}
        \begin{jieda}
            $a_3 + a_6 + a_9 = 3a_6 = 12$，故$a_6 = 4$，因此有$\left \{\begin{array}{l}
                 a_3 + a_9 = 8  \\
                 a_3 \cdot a_9 = 7
            \end{array} \right.$，又因为等差数列$\{a_n\}$为严格增数列，所以$a_3 = 1, a_9 = 7$，故$d = 1$，$a_n = n - 2$.
        \end{jieda}
        \item 若等差数列$\{a_n\}$的公差$d \neq 0$，且$a_1, a_2$为关于$x$的方程$x^2 - a_3x + a_4 = 0$的两根，则此等差数列的通项$a_n = $\blkda{$2n$}
        \begin{jieda}
            由韦达定理有$\left \{ \begin{array}{l}
                 a_1 + a_2 = a_3  \\
                 a_1 \cdot a_2 = a_4 
            \end{array} \right.$，根据第一个式子有$a_1 = d$，代入第二个式子有$a_1 = d = 2$，故$a_n = 2n$.
        \end{jieda}
        \item 若等差数列$\{a_n\}$中$a_2 + a_3 + a_4 + a_5 = 34$，$a_2 a_5 = 52$，且$a_4 > a_2$，则$a_3 = $\blkda{$7$}
        \begin{jieda}
            $a_2 + a_3 + a_4 + a_5 = 2(a_2 + a_5)$，因此有$\left \{\begin{array}{l}
                 a_2 + a_5 = 17  \\
                 a_2 \cdot a_5 = 52
            \end{array} \right.$，结合$a_4 > a_2$可知等差数列$\{a_n\}$为严格增数列，故$a_2 = 4, a_5 = 13$，故公差为$3$，通项公式为$3n-2$，故$a_3 = 7$.
        \end{jieda}
    \end{enumerate}
    \item 根据等差数列性质求值
    \begin{enumerate}
        \item 若$m, a_1, a_2, n$和$m, b_1, b_2, b_3, n$（$m \neq n$）分别是两个等差数列，则$\dfrac{a_2 - a_1}{b_2 - b_1}$的值为\blkda{$\dfrac{4}{3}$}
        \begin{jieda}
            根据第一个等差数列有$m-n = 3(a_2 - a_1)$；根据第二个等差数列有$m-n = 4(b_2 - b_1)$，因此$\dfrac{a_2 - a_1}{b_2 - b_1} = \dfrac{4}{3}$
        \end{jieda}
        \item 若$a, x, b, 2x$依次成等差数列，则$a:b = $\blkda{1:3}
        \begin{jieda}
            $\left \{ \begin{array}{l}
                 a+b = 2x  \\
                 3x = 2b
            \end{array} \right.$，故消去$x$有$3a+3b = 4b$，故$b = 3a$.
        \end{jieda}
        \item 若$a, b, c$的倒数依次成等差数列，且$a, b, c$互不相等，则$\dfrac{a-b}{b-c} = $\blkda{$ \dfrac{a}{c}$}
        \begin{jieda}
            $a, b, c$的倒数依次成等差数列即$\dfrac{1}{c} - \dfrac{1}{b} = \dfrac{1}{b} - \dfrac{1}{a}$，即$\dfrac{b-c}{bc} = \dfrac{a - b}{ab}$，故$\dfrac{a-b}{b-c} = \dfrac{ab}{bc} = \dfrac{a}{c}$
        \end{jieda}
    \end{enumerate}
    \item 等差数列项数的和与项的和的关系
    \begin{enumerate}
        \item 若等差数列$\{a_n\}$中$a_1 + a_3 + a_5 = -1$，则$a_1 + a_2 + a_3 + a_4 + a_5 = $\blkda{$-\dfrac{5}{3}$}
        \begin{jieda}
            $a_1 + a_3 + a_5 = 3a_3$，故$a_3 = -\dfrac{1}{3}$，而$a_1 + a_2 + a_3 + a_4 + a_5 = 5a_3 = -\dfrac{5}{3}$.
        \end{jieda}
        \item 若等差数列$\{a_n\}$中，$S_m = S_n = l$（$m \neq n$），则$a_1 + a_{m+n} = $\blkda{0}
        \begin{jieda}
            不妨设$m<n$,$S_n-S_m = a_{m+1}+\cdots+a_n = 0$
            $a_{m+1} + a_n = 0$，而$a_1 + a_{m+n} = a_{m+1} + a_n$，故$a_1 + a_{m+n} = 0$.
        \end{jieda}
        \item 若等差数列$\{a_n\}$中$a_1 - a_4 - a_8 - a_{12} + a_{15} = 2$，则$a_3 + a_{13} = $\blkda{$-4$}
        \begin{jieda}
            $a_1 - a_4 - a_8 - a_{12} + a_{15} = -a_8$，故$a_8 = -2$，而$a_3 + a_{13} = 2a_8 = -4$
        \end{jieda}
        \item 若等差数列$\{a_n\}$中$a_3 + a_{11} = 10$，则$a_2 + a_4 + a_{15} = $\blkda{15}
        \begin{jieda}
            $a_3 + a_{11} = 2a_7$，故$a_7 = 5$，而$a_2 + a_4 + a_{15} = 3a_7 = 15$.
        \end{jieda}
        \item 若一个等差数列共有$2n+1$项，其奇数项之和为$290$，偶数项之和为$261$，则第$n+1$项为\blkda{$29$}
        \begin{jieda}
            $290 = (n+1)a_{n+1}$，$261 = na_{n+1}$，故$a_{n+1} = 290 - 261 = 29$.
        \end{jieda}
    \end{enumerate}
    \item 等差数列部分和构成等差数列
    \begin{enumerate}
        \item 若等差数列$\{a_n\}$的前4项和为1，前8项和为4，则$a_{17} + a_{18} +a_{19} +a_{20} = $\blkda{9}
        \begin{jieda}
            $a_{1} + a_{2} +a_{3} +a_{4} = 1$；$a_{5} + a_{6} +a_{7} +a_{8} = 3$；$a_{9} + a_{10} +a_{11} +a_{12} = 5$；$a_{13} + a_{14} +a_{15} +a_{16} = 7$；$a_{17} + a_{18} +a_{19} +a_{20} = 9$
        \end{jieda}
        \item 若等差数列$\{a_n\}$的前10项和为100，前20项和为400，则前30项和\blkda{900}
        \begin{jieda}
            $S_{10} = 100, S_{20} - S_{10} = 300 \Rightarrow S_{30} - S_{20} = 500$，故$S_{30} = 900$.
        \end{jieda}
        \item 设等差数列 $\left\{  {a}_{n}\right\}$ 的前 $n$ 项和为 ${S}_{n}$ ,若 ${S}_{3} = 9,{S}_{6} = {36}$ ,则 ${a}_{7} + {a}_{8} + {a}_{9} =$ \blkda{45}.
    \end{enumerate}
    \item 等差数列求和公式应用
    \begin{enumerate}
        \item 若一个等差数列的前10项和是前5项和的4倍，则其首项与公差之比为\blkda{$\dfrac{1}{2}$}
        \begin{jieda}
            $S_{10} = 10a_1 + 45d, S_5 = 5a_1 + 10d$，故$10a_1 + 45d = 4(5a_1 + 10d)$，即$5d = 10a_1$，因此$\dfrac{a_1}{d} = \dfrac{1}{2}$.
        \end{jieda}
        \item 若一个等差数列的前100项和是前10项和的100倍，则$\dfrac{a_{100}}{a_{10}} = $\blkda{$\dfrac{199}{19}$}
        \begin{jieda}
            $S_{100} = 100a_1 + (50\times 99) d, S_{10} = 10a_1 + 45d$，故$100a_1 + (50\times 99) d = 100(10a_1 + 45d)$，即$a_1 + \dfrac{99d}{2} = 10a_1 + 45d$，即$\dfrac{9d}{2} = 9a_1$因此$\dfrac{a_1}{d} = \dfrac{1}{2}$，故$\dfrac{a_{100}}{a_{10}} = \dfrac{199d}{19d} = \dfrac{199}{19}$.
        \end{jieda}
        \item 若等差数列$\{a_n\}$中，$a_1 + a_4 = 10, a_2 - a_3 = 2$，则$\{a_n\}$的前$n$项和$S_n = $\blkda{$-n^2 + 9n$}
        \begin{jieda}
            $a_2 - a_3 = 2$即$d = -2$，故$a_1 + a_4 = 2a_1 + 3d = 10$，故$a_1 = 8$，\\
            因此$S_n = 8n - \dfrac{2n (n-1)}{2} = -n^2 + 9n$.
        \end{jieda}
        \item 若等差数列$\{a_n\}$中，$a_1 + a_3 + a_5 + a_7 + a_9 = \dfrac{25}{2}, a_2 + a_4 + a_6 + a_8 + a_{10} = 15$，则其前$20$项之和$S_{20} = $\blkda{105}
        \begin{jieda}
            $5a_5 = a_1 + a_3 + a_5 + a_7 + a_9 = \dfrac{25}{2}$，即$a_5 = \dfrac{5}{2}$，$5a_6 = a_2 + a_4 + a_6 + a_8 + a_{10} = 15$，即$a_6 = 3$，故$d = \dfrac{1}{2}, a_n = \dfrac{n}{2}$. 故$S_{20} = (\dfrac{1+20}{2}) \times \dfrac{20}{2} = 105$
        \end{jieda}
    \end{enumerate}
    \item 已知等差数列前后$k$项的和
    \begin{enumerate}
        \item 已知 $\left\{  {a}_{n}\right\}$ 为等差数列,且 ${a}_{n - {99}} + {a}_{n - {98}} + \cdots  + {a}_{n - 1} + {a}_{n} = A,{a}_{1} + {a}_{2} + \cdots  + {a}_{100} = B$ , $\left( {n > {100},n \in  {N}^{ * }}\right)$ ,则 ${S}_{n}$ 等于\blkda{$\dfrac{A+B}{200}n$}.
        \item 若等差数列$\{a_n\}$中，前4项的和为21，末4项的和为67，所有项的和为286，则$n = $\blkda{26}
        \begin{jieda}
            $4(a_1 + a_n) = 21 + 67$，故$(a_1 + a_n) = 22$，故$22 \times \dfrac{n}{2} = 286$，解得$n = 26$
        \end{jieda}
        \item 若等差数列$\{a_n\}$中，前3项的和为12，末3项的和为75，各项之和为145，则项数$n = $\blkda{10}，公差$d = $\blkda{3}，首项$a_1 = $\blkda{1}
        \begin{jieda}
            $3(a_1 + a_n) = 12 + 75$，故$(a_1 + a_n) = 29$，故$29 \times \dfrac{n}{2} = 145$，解得$n = 10$. \\
            由前3项的和为12，末3项的和为75知$a_2 = 4, a_9 = 25$，故$d = 3$，因此$a_1 = a_2 - 3 = 1$.
        \end{jieda}
    \end{enumerate}
    \item 等比数列的定义
    \begin{enumerate}
        \item 等比数列$\{a_n\}$的公比为$q$，$a_n = a_m \cdot x$，则$x = $\blkda{$q^{n-m}$}.
        \item 设数列$\{a_n\}$时各项为正数的无穷数列，$A_i$是边长为$a_i, a_{i+1}$的矩形面积（$i \in \mathbf{N^*}$），则数列$\{A_n\}$是等比数列的充要条件是\dotfill{(\kh{D})}
        \xx{数列$\{a_n\}$是等比数列}{$a_1, a_3, \cdots a_{2n-1}$或$a_2, a_4, \cdots a_{2n}$是等比数列}{$a_1, a_3, \cdots a_{2n-1}$和$a_2, a_4, \cdots a_{2n}$是等比数列}{$a_1, a_3, \cdots a_{2n-1}$和$a_2, a_4, \cdots a_{2n}$是具有相同公比的等比数列}
        \item 已知$a, b, c, x, y, z$都是不为1的正数，且有$a^x = b^y = c^z$，如果$\dfrac{1}{x}, \dfrac{1}{y}, \dfrac{1}{z}$成等差数列，求证$a, b, c$成等比数列.
        \begin{jieda}
            设$k = a^x = b^y = c^z$，则$x = log_ak, y = log_bk, z = log_ck$，故$\dfrac{1}{x} = log_ka, \dfrac{1}{y} = log_kb, \dfrac{1}{z} = log_kc$. $\dfrac{1}{x}, \dfrac{1}{y}, \dfrac{1}{z}$成等差数列即$log_k(ac) = log_k(b^2)$，故$\dfrac{a}{b} = \dfrac{b}{c}$，因此$a, b, c$成等比数列.
        \end{jieda}
    \end{enumerate}
    \item 求等比数列的通项公式、项和公比
    \begin{enumerate}
        \item 在等比数列$\{a_n\}$中，若$a_5 = 2, a_{10} = 10$，则$a_{15} = $\blkda{50}.
        \item 在等比数列$\{a_n\}$中，$a_2 = 5, a_4 = \dfrac{1}{5}$，则$a_1 = $\blkda{25或-25}
        \begin{jieda}
            $\dfrac{a_4}{a_2} = q^2$，解得$q = \pm \dfrac{1}{5}$，因此$a_1 = \dfrac{a_2}{q} = $25或-25.
        \end{jieda}
        \item 在等比数列$\{a_n\}$中，若$a_1 + a_2 + a_3 = -3, a_1 a_2 a_3 = 8$，则$a_4 = $\blkda{8或$\dfrac{1}{2}$}
        \begin{jieda}
            $a_1 a_2 a_3 = 8 \Rightarrow a_2 = 2$，故有$\left \{ \begin{array}{l}
                 a_1 + a_3 = -5  \\
                 a_1 \cdot a_3 = 4
            \end{array} \right .$，解得$\left \{\begin{array}{l}
                 a_1 = -1  \\
                 a_3 = -4
            \end{array} \right.$或$\left \{\begin{array}{l}
                 a_1 = -4  \\
                 a_3 = -1
            \end{array} \right.$，故公比$q = -2$或$-\dfrac{1}{2}$，$a_4 = $8或$\dfrac{1}{2}$
        \end{jieda}
        \item 在等比数列$\{a_n\}$中，$a_1 + a_2 + a_3 = 26$，$a_4 - a_1 = 52$，则$a_n = $\blkda{$2\cdot 3^{n-1}$}
        \begin{jieda}
            $a_1 + a_2 + a_3 = 26$即$a_1 \cdot (1 + q + q^2) = 26$，$a_4 - a_1 = 52$即$a_1(q^3 - 1) = a_1 (q-1)(1 + q + q^2) = 52$，故$q = 3$，因此$a_n = 2\cdot 3^{n-1}$.
        \end{jieda}
        \item 在等比数列$\{a_n\}$中，$a_1 + a_6 = 11, a_3 \cdot a_4 = \dfrac{32}{9}$，且$\dfrac{2}{3}a_2, a_3^2, a_4+\dfrac{4}{9}$依次是等差数列，则$a_n = $\blkda{$\dfrac{2^{n-1}}{3}$}
        \begin{jieda}
            $a_3 \cdot a_4 = a_1 \cdot a_6$，故$\left \{ \begin{array}{l}
                 a_1 + a_6 = 11  \\
                 a_1 \cdot a_6 = \dfrac{32}{9}
            \end{array} \right .$，解得$\left \{\begin{array}{l}
                 a_1 = \dfrac{1}{3}  \\
                 a_6 = \dfrac{32}{3}
            \end{array} \right.$或$\left \{\begin{array}{l}
                 a_1 = \dfrac{32}{3}  \\
                 a_6 = \dfrac{1}{3}
            \end{array} \right.$，故$a_1 = \dfrac{1}{3}, q = 2$或$a_1 = \dfrac{32}{3}, q = \dfrac{1}{2}$，代入“$\dfrac{2}{3}a_2, a_3^2, a_4+\dfrac{4}{9}$依次是等差数列”验证可知仅有$a_1 = \dfrac{1}{3}, q = 2$符合，因此$a_n = \dfrac{2^{n-1}}{3}$.
        \end{jieda}
        
        \item 已知等比数列$\{a_n\}$的公比$q>1$，第19项的平方等于第26项，则使得$a_1 + a_2 + \cdots + a_n > \dfrac{1}{a_1} + \dfrac{1}{a_2} + \cdots + \dfrac{1}{a_n}$成立的$n$的范围是\blkda[3]{$\{n | n \geq 24, n \in \mathbf{N}\}$}
        \begin{jieda}
            因为$q > 1$，所以$\{a_n\}$是递增数列，故当$n$大于某值后该不等式必然恒成立. \\
            $(a_1 \cdot q^{18})^2 = (a_1 \cdot q^{25})$，故$a_1 = \dfrac{1}{q^{11}}$，故$a_{12} = 1$，故$a_1 + a_2 + \cdots + a_{23} = \dfrac{1}{a_1} + \dfrac{1}{a_2} + \cdots + \dfrac{1}{a_{23}}$，因此当$n \geq 24$时不等式恒成立.
        \end{jieda}
    \end{enumerate}
    \item 等比中项 \& 等比数列的部分和构成等比数列 
    \begin{enumerate}
    \item $m,n,p,q \in \mathbf{N^*}$，且有$m+n = p+r$则等比数列$\{a_n\}$必定满足\dotfill{(\kh{A})}
    \xx{$\dfrac{a_m}{a_p} = \dfrac{a_r}{a_n}$}
    {$\dfrac{a_m}{a_n} = \dfrac{a_r}{a_p}$}
    {$a_m + a_n = a_p + a_r$}
    {$a_m- a_n = a_p- a_r$}
    \item 在等比数列$\{a_n\}$中，若$a_4 = 5, a_8 = 6$，则$a_2 a_{10} = $\blkda{30}
    \begin{jieda}
        $a_2 a_{10} = a_4 a_8 = 30$.
    \end{jieda}
    \item 已知$a = 243, b = 3$，$c$是$a, b$的等比中项
    \begin{enumerate}
        \item 若$d, e$分别是$a, c$和$c, b$的等差中项，则$a+b+c+d+e = $\blkda[4]{423或315}
        \item 若$d, e$分别是$a, c$和$c, b$的等比中项，则$a+b+c+d+e = $\blkda[4]{363或345或201或183}
    \end{enumerate}
    \begin{jieda}
        首先可以得到$c = -27$或27.
        \begin{enumerate}
            \item $a+b+c+d+e = a+b+c + \dfrac{a+c}{2} + \dfrac{b+c}{2} = \dfrac{3(a+b)}{2} + 2c = $ 423或315
            \item 此时必有$c = 27$，\begin{enumerate}[1.]
                \item $d = 81, e = 9$，$a+b+c+d+e = 363$
                \item $d = 81, e = -9$，$a+b+c+d+e = 345$
                \item $d = -81, e = 9$，$a+b+c+d+e = 201$
                \item $d = -81, e = -9$，$a+b+c+d+e = 183$
            \end{enumerate}
        \end{enumerate}
    \end{jieda}
    \item 已知正项等比数列$\{a_n\}$的公比不是1，则\dotfill{(\kh{A})}
    \xx{$a_2 + a_8 > a_3 + a_7$}{$a_2 + a_8 < a_3 + a_7$}{仅当$q>1$时，$a_2 + a_8 > a_3 + a_7$}{仅当$1>q>0$时，$a_2 + a_8 > a_3 + a_7$}
    \begin{jieda}
        $(a_2 + a_8) - (a_3 + a_7) = (a_2-a_3)+(a_8-a_7) = a_2(1-q)+a_7(q-1) = (a_2-a_7)(1-q) = a_2(1-q^5)(1-q)$
            
        易知$(1-q^5)(1-q) > 0$在$q > 0$且$q \neq 1$时恒成立，因此$a_2 + a_8 > a_3 + a_7$恒成立.
    \end{jieda}
    \item 若等比数列$\{a_n\}$的公比$q < 0$，前$n$项和为$S_n$则$S_8a_9$与$a_8S_9$的大小关系是\dotfill{(\kh{A})}
    \xx{$S_8a_9 > a_8S_9$}{$S_8a_9 < a_8S_9$}{$S_8a_9 = a_8S_9$}{无法确定}
    \begin{jieda}
        $a_8S_9 = a_8(a_1 + qS_8) = a_1a_8 + S_8a_9$\\
        因为等比数列$\{a_n\}$的公比$q < 0$，所以$a_1a_8 < 0$，故$S_8a_9 > a_8S_9$
    \end{jieda}
    \end{enumerate}
    \item 等差数列与等比数列综合
    \begin{enumerate}
        \item 若互不相等的实数$a, b, c$成等差数列，$c, a, b$成等比数列，且有$a+3b+c = 10$，则$a = $\blkda{-4}
        \begin{jieda}
            $\left \{ \begin{array}{ll}
                 a+c = 2b & \mbox{\ding{172}} \\
                 a+3b+c = 10 & \mbox{\ding{173}} \\
                 b \cdot c = a^2& \mbox{\ding{174}}
            \end{array}\right .$，首先将\ding{172}代入\ding{173}消去$b$得\ding{175} $a+c = 4$，将\ding{172}代入\ding{174}得\ding{176}$ac + c^2 = 2a^2$，将\ding{175}代入\ding{176}后得$a(4-a) + (4-a)^2 = 2a^2$，解得$a = 2$或$-4$.当$a = 2$时$c = 2$与互不相等的要求矛盾，故$a = -4$.
        \end{jieda}
        \item 在等差数列$\{a_n\}$中，$a_1, a_3, a_9$成等比数列，且公差$d$不为0，则$\dfrac{a_1 + a_3 + a_9}{a_2 + a_4 + a_{10}} = $\blkda{$\dfrac{13}{16}$}.
        \begin{jieda}
            易知$a_n = n \cdot a_1$，故$\dfrac{a_1 + a_3 + a_9}{a_2 + a_4 + a_{10}} = \dfrac{1+3+9}{2+4+10} = \dfrac{13}{16}$，下面为严格推导：\\
            $a_1 \cdot (a_1+8d) = (a_1 + 2d)^2$，即$a_1^2 + 8a_1d = a_1^2 + 4a_1d + 4d^2$，即$a_1 = d$，故$a_n = n \cdot a_1$
        \end{jieda}
        \item 数列$\{a_n\}$既是等差数列又是等比数列，且有$a_{n+1} = \dfrac{a_n^2}{2a_n - 1}$，则数列$\{a_n\}$的前10项和为\blkda{10}
        \begin{jieda}
            数列$\{a_n\}$既是等差数列又是等比数列即$\{a_n\}$为非零常数列，设$a_n = k$，则$k = \dfrac{k^2}{2k-1}$，故$k = 1$，因此$a_n = 1$，前10项和为10.
        \end{jieda}
    \end{enumerate}
    \item 构造新数列求通项
    \begin{enumerate}
        \item $a_{n+1} = pa_n + q (p \neq 1)$型 \\
            若数列$\{a_n\}$满足$a_1 = 1, a_{n+1} = 3a_n+1$，则$a_n = $\blkda{$\dfrac{1}{2}(3^n - 1)$}
            \begin{jieda}
                设$a_{n+1} + x = 3(a_n + x)$，解得$x = \dfrac{1}{2}$. 因此令$b_n = a_n + \dfrac{1}{2}$，则$b_1 = \dfrac{3}{2}, b_{n+1} = 3b_n$，因此$\{b_n\}$是等比数列，$b_n = \dfrac{1}{2} \cdot 3^n$，故$a_n = \dfrac{1}{2}(3^n - 1)$ 
            \end{jieda}
        \item $a_{n+1} = pa_n + An+B (p \neq 1)$型 \\
            若数列$\{a_n\}$满足$a_1 = 1, a_{n+1} = 3a_n + n + 1$，则$a_n = $\blkda[2.5]{$\dfrac{3^{n+1}-2n-3}{4}$}
            \begin{jieda}
                设$a_{n+1} + x \cdot (n+1) + y = 3(a_n + x \cdot n + y)$，则$\left \{\begin{array}{l}
                    3x - x = 1  \\
                    3y - y - x = 1
                \end{array} \right.$，解得$x = \dfrac{1}{2}, y = \dfrac{3}{4}$.因此令$b_n = a_n + \dfrac{n}{2} + \dfrac{3}{4}$，则$b_1 = \dfrac{9}{4}, b_{n+1} = 3b_n$，因此$\{b_n\}$是等比数列，$b_n = \dfrac{3}{4} \cdot 3^n$，故$a_n = \dfrac{3}{4} \cdot 3^n - \dfrac{n}{2} - \dfrac{3}{4} = \dfrac{3^{n+1}-2n-3}{4}$
            \end{jieda}
        \item $a_{n+1} = pa_n + An^2+Bn+C (p \neq 1)$ 型 \\
        若数列$\{a_n\}$满足$a_1 = 1, a_{n+1} = 3a_n + 2n^2 + n + 1$，则$a_n = $\blkda[3.5]{$\dfrac{7 \cdot 3^n - 4n^2 - 6n - 7}{4}$}
        \begin{jieda}
            设$a_{n+1} + x \cdot (n+1)^2 + y \cdot (n+1) + z = 3(a_n + x \cdot n^2 + y \cdot n + z)$，则$\left \{ \begin{array}{l}
                 3x - x = 2  \\
                 3y - y - 2x = 1  \\
                 3z - z - y - x = 1
            \end{array} \right.$，解得$\left \{ \begin{array}{l}
                 x = 1  \\
                 y = \dfrac{3}{2}  \\
                 z = \dfrac{7}{4}
            \end{array} \right.$. 因此令$b_n = a_n + n^2 + \dfrac{3n}{2} + \dfrac{7}{4}$，则$b_1 = \dfrac{21}{4}, b_{n+1} = 3b_n$，因此$\{b_n\}$是等比数列，$b_n = \dfrac{7}{4} \cdot 3^n$，故$a_n = \dfrac{7}{4} \cdot 3^n - n^2 - \dfrac{3n}{2} - \dfrac{7}{4} = \dfrac{7 \cdot 3^n - 4n^2 - 6n - 7}{4}$
        \end{jieda}
        \item $a_{n+1} = pa_n + p^{n+1} (p \neq 1)$型 \\
        若数列$\{a_n\}$满足$a_1 = 1, a_{n+1} = 2a_n + 2^{n+1}$，则$a_n = $\blkda[2.5]{$(2n-1)\cdot 2^{n-1}$}
        \begin{jieda}
            $a_{n+1} = 2a_n + 2^{n+1}$即$\dfrac{a_{n+1}}{2^{n+1}} = \dfrac{a_n}{2^n} + 1$，设$b_n = \dfrac{a_n}{2^n}$，则$b_1 = \dfrac{1}{2}, b_{n+1} = b_n + 1$，因此$\{b_n\}$是等差数列，$b_n = \dfrac{2n-1}{2}$，故$a_n = (2n-1)\cdot 2^{n-1}$
        \end{jieda}
        \item $a_{n+1} = pa_n + q^{n+1} (p \neq 1)$型 \\
        若数列$\{a_n\}$满足$a_1 = 1, a_{n+1} = 2a_n + 3^{n+1}$，则$a_n = $\blkda[2.5]{$3^{n+1} - 2^{n+2}$}
        \begin{jieda}
            $a_{n+1} = 2a_n + 3^{n+1}$即$\dfrac{a_{n+1}}{3^{n+1}} = \dfrac{2}{3} \cdot \dfrac{a_n}{3^n} + 1$，设$b_n = \dfrac{a_n}{3^n}$，则$b_1 = \dfrac{1}{3}, b_{n+1} = \dfrac{2}{3}b_n + 1$，设$c_n = b_n - 3$，则$c_1 = -\dfrac{8}{3}, c_{n+1} = \dfrac{2}{3}c_n$，因此$\{c_n\}$是等比数列，$c_n = -4 \cdot \left ( \dfrac{2}{3} \right )^n$，故$b_n = 3 - 4 \cdot \left ( \dfrac{2}{3} \right )^n$，$a_n = 3^{n+1} - 2^{n+2}$
        \end{jieda}
        \item $a_{n+1} = pa_n + k \cdot q^{n+1} (p \neq 1)$型 \\
        若数列$\{a_n\}$满足$a_1 = 1, a_{n+1} = 2a_n + 4 \cdot 3^{n+1}$，则$a_n = $\blkda[3.5]{$4 \cdot 3^{n+1} - 35 \cdot 2^{n-1}$}
        \begin{jieda}
            $a_{n+1} = 2a_n + 4 \cdot 3^{n+1}$即$\dfrac{a_{n+1}}{3^{n+1}} = \dfrac{2}{3} \cdot \dfrac{a_n}{3^n} + 4$，设$b_n = \dfrac{a_n}{3^n}$，则$b_1 = \dfrac{1}{3}, b_{n+1} = \dfrac{2}{3}b_n + 4$，设$c_n = b_n - 12$，则$c_1 = -\dfrac{35}{3}, c_{n+1} = \dfrac{2}{3}c_n$，因此$\{c_n\}$是等比数列，$c_n = -\dfrac{35}{2} \cdot \left ( \dfrac{2}{3} \right )^n$，故$b_n = 12 - \dfrac{35}{2}  \cdot \left ( \dfrac{2}{3} \right )^n$，$a_n = 4 \cdot 3^{n+1} - 35 \cdot 2^{n-1}$
        \end{jieda}
        \item $a_{n+2} = pa_{n+1} + qa_n$型 \\
        若数列$\{a_n\}$的首项$a_1 = 1, a_2 = 4$，且满足$a_{n+2} = 2a_{n+1}+3a_n$，则$a_n = $\blkda[3]{$\dfrac{5 \cdot 3^{n-1} + (-1)^n}{4}$}
        \begin{jieda}
            设$a_{n+2} + \lambda a_{n+1} = \mu(a_{n+1} + \lambda a_n)$，则$\left \{ \begin{array}{l}
                \mu - \lambda = 2 \\
                \mu \cdot \lambda = 3
            \end{array} \right.$，取$\left \{ \begin{array}{l}
                \lambda = 1 \\
                \mu = 3
            \end{array} \right.$，则有$a_{n+2} + a_{n+1} = 3(a_{n+1} + a_n)$，设$b_n = a_n + a_{n+1}$，则$b_1 = 5, b_{n+1} = 3b_n$，因此$\{b_n\}$是等比数列，$b_n = 5 \cdot 3^{n-1}$，即$a_{n+1} = -a_n + 5 \cdot 3^{n-1}$，故$\dfrac{a_{n+1}}{3^{n+1}} = -\dfrac{1}{3} \cdot \dfrac{a_n}{3^n} + \dfrac{5}{9}$，设$c_n = \dfrac{a_n}{3^n}$，则$c_1 = \dfrac{1}{3}, c_{n+1} = -\dfrac{1}{3} c_n + \dfrac{5}{9}$，设$d_n = c_n - \dfrac{5}{12}$，则$d_1 = -\dfrac{1}{12}, d_{n+1} = -\dfrac{1}{3}\cdot d_n$，故$\{d_n\}$是等比数列，$d_n = \dfrac{1}{4} \cdot \left (-\dfrac{1}{3} \right)^n$，故$c_n = \dfrac{5}{12} + \dfrac{1}{4} \cdot \left (-\dfrac{1}{3} \right)^n$，$a_n = \dfrac{5 \cdot 3^{n-1} + (-1)^n}{4}$
        \end{jieda}
        \item $a_{n+1} = \dfrac{pa_n}{qa_n+r}$型 \\
        若数列$\{a_n\}$满足$a_1 = 1, a_{n+1} = \dfrac{a_n}{2a_n + 3}$，则$a_n = $\blkda[2.5]{$\dfrac{1}{2\cdot 3^{n-1}-1}$}
        \begin{jieda}
            $a_{n+1} = \dfrac{a_n}{2a_n + 3}$即$\dfrac{1}{a_{n+1}} = 2 + \dfrac{3}{a_n}$，设$b_n = \dfrac{1}{a_n}$，则$b_1 = 1, b_{n+1} = 3b_n + 2$，设$c_n = b_n + 1$，则$c_1 = 2, c_{n+1} = 3c_n$，故$\{c_n\}$是等比数列，$c_n = 2 \cdot 3^{n-1}$，故$b_n = 2\cdot 3^{n-1}-1, a_n = \dfrac{1}{2\cdot 3^{n-1}-1}$.
        \end{jieda}
        \item $a_{n+1} = pa_n^k$型 \\
        若数列$\{a_n\}$满足$a_1 = 1, a_{n+1} = 2a_n^2$，则$a_n = $\blkda{$2^{2^{n-1}-1}$}
        \begin{jieda}
            $a_{n+1} = 2a_n^2$即$log_2(a_{n+1}) = 2 \cdot log_2(a_n) + 1$，设$b_n = log_2(a_n)$，则$b_1 = 0, b_{n+1} = 2b_n + 1$，故设$c_n = b_n + 1$，则$c_1 = 1, c_{n+1} = 2c_n$，故$c_n = 2^{n-1}$，故$b_n = 2^{n-1}-1$，因此$a_n = 2^{2^{n-1}-1}$.
        \end{jieda}
    \end{enumerate}

    \item 分式型递推关系转化为等差数列、等比数列的递推关系
    \begin{enumerate}
        \item 若数列$\{a_n\}$满足$a_1 = 1$，$\dfrac{a_n}{a_n + a_{n+1}} = 2(n \in \mathbf{N^*})$，则$a_n = $\blkda{$\left ( -\dfrac{1}{2}\right)^{n-1}$}
        \begin{jieda}
            $\dfrac{a_n}{a_n + a_{n+1}} = 2$即$a_n = 2a_n + 2a_{n+1}$，即$a_{n+1} = -\dfrac{1}{2}a_n$，故$a_n = \left ( -\dfrac{1}{2}\right)^{n-1}$
        \end{jieda}
        \item 已知数列$\{a_n\}$各项均不为0，且$a_1 = 3$，$\dfrac{a_n}{a_{n+1} - a_n} - \dfrac{a_{n-1}}{a_n - a_{n-1}} = 1 (n \geq 2 \mbox{且} n \in \mathbf{N})$，若$a_2 = 8$，则$a_n = $\blkda{$5n-2$}
        \begin{jieda}
            $\dfrac{a_n}{a_{n+1} - a_n} - \dfrac{a_{n-1}}{a_n - a_{n-1}} = 1 \Rightarrow \dfrac{a_n}{a_{n+1} - a_n}  =  \dfrac{a_n}{a_n - a_{n-1}} \Rightarrow a_{n+1} - a_n = a_n - a_{n-1}$，故$\{a_n\}$是等差数列，$d = 8 - 3 = 5$，因此$a_n = 5n-2$
        \end{jieda}
    \end{enumerate}
    \item 由递推关系推理数列性质
    \begin{enumerate}
        \item 已知数列$\{a_n\}$中，$a_1 = 3, a_{n+1} = \sqrt{a_n + 2}$ \dotfill{(\kh{A})}
        \xx{数列$\{a_n\}$严格减}{数列$\{a_n\}$严格增}{数列$\{a_n\}$先减后增}{数列$\{a_n\}$先增后减}
        \begin{jieda}
            $y = x$与$y = \sqrt{x+2}$相交于$(2,2)$，当$x > 2$时$y = x$在$y = \sqrt{x+2}$上方. 因为$a_1 = 3$，所以$\{a_n\}$为严格减数列，趋于$2$.
        \end{jieda}
        \item 已知数列$\{a_n\}$各项都是正数，且$a^2_{n+1} - a_n = d$（$d$为常数，$n \in \mathbf{N^*}$），给出下列四个结论：
        \begin{dingautolist}{172}
            \item $\{a_n\}$的前$n$项和$S_n$有最小值.
            \item 若$\{a_n\}$为严格增数列，则$d>0$
            \item 若$\{a_n\}$为周期数列，则最小正周期可能为$2$.
            \item 若$\{a_n\}$为常数列，则$d \geq -\dfrac{1}{4}$
        \end{dingautolist}其中正确的结论是\blkda{\ding{172} \ding{173} \ding{175}}
        \begin{jieda}
            \begin{dingautolist}{172}
                \item $\{a_n\}$各项均为正数，故必有$\{S_n\}$严格增，故其最小值为$S_1 = a_1$.
                \item $a^2_{n+1} - a_n = d$，因为$\{a_n\}$各项均为正数，所以$a_{n+1} = \sqrt{a_n + d}$，故$\{a_n\}$为严格增数列的必要条件是$y = \sqrt{x+d}$有部分图像在$y = x$上方，即$d \geq -\dfrac{1}{4}$，故可举出反例：$d = 0, a_1 = \dfrac{1}{2}$
                \item 根据$y = \sqrt{x+d}$与$y = x$的关系可知$\{a_n\}$除常数列之外，不可能是周期数列
                \item 只要$y = \sqrt{x+d}$与$y = x$有交点，$\{a_n\}$就可常数列.
            \end{dingautolist}
        \end{jieda}
        \item 已知数列$\{a_n\}$满足$a_n > 0$，$a_{n+1}^2 - a_{n+1} = a_n$，给出下列四个结论：
        \begin{dingautolist}{172}
            \item 对于任意$n \geq 3$，都有$a_n \geq 2$
            \item $\{a_n\}$不可能为常数列
            \item 若$0 < a_1 < 2$，则$\{a_n\}$为递增数列
            \item 若$a_1 > 2$，则$2 < a_n < a_1 (n \geq 2)$
        \end{dingautolist}其中正确的结论是\blkda{\ding{174} \ding{175}}
        \begin{jieda}
            $y = \dfrac{1+\sqrt{1+4x}}{2} (x > 0)$（即抛物线$y^2 - y = x (x> 0, y > 1)$）与$y = x$交于$(2, 2)$，故$a_1 \in (0, 2)$时，$\{a_n\}$为严格增数列，趋于2；$a_1 \in (2, +\infty)$时，$\{a_n\}$为严格减数列，趋于2；而$a_1 = 2$时$\{a_n\}$为常数列.
        \end{jieda}
        \item 已知数列$\{a_n\}$满足$a_{n+1} = \dfrac{1}{4}(a_n - 6)^3 + 6$，则\dotfill{(\kh{B})}
        \xx{当$a_1 = 3$时，$\{a_n\}$为递减数列，且存在常数$M \leq 0$，使得$a_n > M$恒成立}
        {当$a_1 = 5$时，$\{a_n\}$为递增数列，且存在常数$M \leq 6$，使得$a_n < M$恒成立}
        {当$a_1 = 7$时，$\{a_n\}$为递减数列，且存在常数$M > 6$，使得$a_n > M$恒成立}
        {当$a_1 = 9$时，$\{a_n\}$为递增数列，且存在常数$M > 0$，使得$a_n < M$恒成立}
        \begin{jieda}
            设$b_n = a_n - 6$，则$b_{n+1} = \dfrac{b_n^3}{4}$，$y = x$与$y = \dfrac{x^3}{4}$交于$(0,0),(2, 2), (-2, -2)$，根据$y = x$与$y = \dfrac{x^3}{4}$的关系可知：
            \begin{dingautolist}{172}
                \item $b_1 \in (-\infty, -2)$，此时$\{b_n\}$为递减数列，且趋于$-\infty$
                \item $b_1 \in (-2, 0)$，此时$\{b_n\}$为递增数列，且趋于$0$
                \item $b_1 \in (0, 2)$，此时$\{b_n\}$为递减数列，且趋于$0$
                \item $b_1 \in (2, +\infty)$，此时$\{b_n\}$为递增数列，且趋于$+\infty$
            \end{dingautolist}
            因此选B，对应的$M = 6$.
        \end{jieda}

        \item 已知数列 $\left\{  {a}_{n}\right\}$ 满足 ${a}_{n + 1} = \left| {{a}_{n} + 1}\right|  + \lambda \left| {{a}_{n} - 1}\right|$ ，其中 $\lambda$ 为常数. 对于下述两个命题:

        \ding{172}  对于任意的 $\lambda  > 0$ ，任意的 ${a}_{1} \in  \mathbf{R}$ ，都有 $\left\{  {a}_{n}\right\}$ 是严格增数列；

        \ding{173}  对于任意的 $\lambda  < 0$ ，存在 ${a}_{1} \in  \mathbf{R}$ ，使得 $\left\{  {a}_{n}\right\}$ 是严格减数列.

        以下说法正确的为 \dotfill{(\kh{A})}

        \xx{\ding{172} 真命题， \ding{173} 假命题}{\ding{172} 假命题， \ding{173} 真命题}{\ding{172}  \ding{173} 均为真命题}{\ding{172}  \ding{173} 均为假命题} 
        \begin{jieda}
            \begin{dingautolist}{172}
                \item 当$\lambda > 0$时，易知$a_n > 0, n > 1$，而$y = |x+1|+\lambda|x-1|$没有不动点，根据图象可知严增.
                \item 当$\lambda = -1$时，根据$y = |x+1|-|x-1|$的图像结合不动点分类讨论可知命题为假.
            \end{dingautolist}
        \end{jieda}

        \item 已知数列 $\left\{  {a}_{n}\right\}$ 满足 ${a}_{n + 1} = r{a}_{n}\left( {1 - {a}_{n}}\right) \left( {n = 1,2,3,\cdots }\right)$ ， ${a}_{1} \in  \left( {0,1}\right)$ ，给出以下四个结论:

        \ding{172}  当 $r = 2$ 时，只存在有限个 ${a}_{1}$ ，使得对任意正整数 $n$ ，都有 ${a}_{n + 1} > {a}_{n}$ ；

        \ding{173}  当 $r = 2$ 时，存在 ${a}_{1}$ 和正整数 $P$ ，当 $n > P$ 时， ${a}_{n + 1} - {a}_{n} < \dfrac{1}{2025}$ ；

        \ding{174}  当 $r = 3$ 时，存在 ${a}_{1}$ 和正整数 $P$ ，当 $n > P$ 时， ${a}_{n + 1} = {a}_{n}$ ；

        \ding{175}  当 $r =  - 3$ 时,不存在 ${a}_{1}$ ,使得对任意正整数 $n$ ,且 $n \geq  3$ ,都有 ${a}_{n} > 0$. 
        
        其中正确结论是 \dotfill{(\kh{B})}

        \xx{\ding{172}  \ding{173}}{\ding{173}  \ding{174}}{\ding{174}  \ding{175}}{\ding{173}  \ding{175}}
        \begin{jieda}
            数形结合，结合不动点判断即可.
        \end{jieda}
    \end{enumerate}
    \item 累加法
    \begin{enumerate}
        \item 若数列$\{a_n\}$满足$a_1 = 1, a_{n+1} = a_n + 2^n$，则$a_{10} = $\blkda{1023}
        \begin{jieda}
            设$b_n = 2^n$，则数列$\{b_n\}$的前$n$项和为$S_n = 2^{n+1} - 2$，由$a_{n+1} = a_n + b_n$可知根据累加法有$a_n = a_1 + S_{n-1} = 2^n - 1(n \geq 2)$，故$a_{10} = 1023$
        \end{jieda}
        \item 若数列$\{a_n\}$满足$a_1 = 2, a_{n+1} = a_n + 3 \cdot 2^n$，则$a_n = $\blkda{$3 \cdot 2^n - 4$}
        \begin{jieda}
            根据$a_{n+1} = a_n + 3 \cdot 2^n$累加得$a_n = 2 + 3(2 + 4 + \cdots + 2^{n-1}) = 3 \cdot 2^n - 4(n \geq 2)$，经检验$a_n = 3 \cdot 2^n - 4$.
        \end{jieda}
        \item 若数列$\{a_n\}$满足$a_1 = a_2 = 1, a_n + a_{n+1} + a_{n+2}= n^2 (n \in \mathbf{N^*})$，则$a_{100} = $\blkda{3268}
        \begin{jieda}
            $\left \{ \begin{array}{l}
                a_n + a_{n+1} + a_{n+2}= n^2  \\
                a_{n+1} + a_{n+2} + a_{n+3}= (n+1)^2 
            \end{array}\right.$，因此$a_{n+3} - a_n = 2n+1$，故$a_{100} = a_1 + 2 \times \sum_{i = 1}^{33}(3i-2) + 33 = 3268$.
        \end{jieda}
        \item 已知数列$\{a_n\}, a_1 = 1$，对任意正整数$k$，$a_{2k-1}, a_{2k}, a_{2k+1}$成等差数列，公差为$2k$，则$a_{100} = $\blkda{5001}
        \begin{jieda}
            $\left \{ \begin{array}{l}
                 a_{2k} - a_{2k-1} = 2k  \\
                 a_{2k-1} - a_{2k-2} = 2k-2 
            \end{array}\right .$，因此$a_k - a_{2k-2} = 4k-2$，故累加得$a_{100} = a_2 + (198 + 194 + \cdots + 6) = a_2 + 4998 = a_1 + 2 + 4998 = 5001$.
        \end{jieda}
        % \item 已知数列$\{a_n\}$满足$a_1 = 1$，$a_{n+1} = a_n - \dfrac{1}{3}a_n^2(n \in \mathbf{N^*})$，则\dotfill{(\kh{B})}
        % \xx{$2 < 100a_{100} < \dfrac{5}{2}$}{$\dfrac{5}{2} < 100a_{100} < 3$}{$3 < 100a_{100} < \dfrac{7}{2}$}{$\dfrac{7}{2} < 100a_{100} < 4$}
        % \begin{jieda}
        %     本题的目标是推理$a_n$的取值范围，首先根据递推关系可以确定$a_n \in (0, 1) (n \geq 2)$，且数列为严格减数列. 为了推理出$a_n$的取值范围，应考虑采用累加法，但直接移项进行累加并不能得到好的性质，因此考虑倒数法：\\
        %     $a_{n+1} = a_n(1 - \dfrac{a_n}{3})$，$\dfrac{1}{a_{n+1}} = \dfrac{1}{a_n \cdot (1 - \dfrac{a_n}{3})} = \dfrac{3}{a_n \cdot (3 - a_n)} = \dfrac{1}{a_n} + \dfrac{1}{3-a_n}$，因此$\dfrac{1}{a_{n+1}} - \dfrac{1}{a_n} = \dfrac{1}{3-a_n} > \dfrac{1}{3}$，故累加可得$\dfrac{1}{a_n} > \dfrac{n-1}{3} + \dfrac{1}{a_1} = \dfrac{n+2}{3}$，即$a_n < \dfrac{3}{n+2}$（$n \geq 2$），故$a_n \leq \dfrac{3}{n+2}$（当且仅当$n = 1$时等号成立）. 代入$\dfrac{1}{a_{n+1}} - \dfrac{1}{a_n} = \dfrac{1}{3-a_n}$有$\dfrac{1}{a_{n+1}} - \dfrac{1}{a_n} < \dfrac{1}{3-\dfrac{3}{n+2}} = \dfrac{n+2}{3n+3} = \dfrac{1}{3}(1+\dfrac{1}{n+1})$，故累加可得$\dfrac{1}{a_n} - 1 < \dfrac{n-1}{3} + \dfrac{1}{3}(\dfrac{1}{2} + \dfrac{1}{3} + \cdots + \dfrac{1}{n})$ \\
        %     综上，$100a_{100} < \dfrac{300}{102} < 3$，\\
        %     $\dfrac{1}{a_{100}} < 1 + 33 + \dfrac{1}{3}(\dfrac{1}{2} + \dfrac{1}{3} + \cdots + \dfrac{1}{100}) < 34 + \dfrac{1}{3}(\dfrac{1}{2} \times 4 + \dfrac{1}{6} \times 96) = 40$，故$100a_{100} > \dfrac{5}{2}$
        % \end{jieda}
    \end{enumerate}
    \item 累乘法
    \begin{enumerate}
        \item 若数列$\{a_n\}$满足$(n-1)a_n = (n+1)a_{n-1}(n \geq 2), a_1 = 2$，则满足$a_n < 930$的最大正整数$n = $\blkda{29}
        \begin{jieda}
            根据递推关系可知$a_n \neq 0$，故$\dfrac{a_n}{a_{n-1}} = \dfrac{n+1}{n-1}$，故累乘得到$a_n = n(n+1)(n \geq 2)$，经检验$a_n = n(n+1)$. 故$a_{30} = 930$，因此$n = 29$.
        \end{jieda}
        \item 若数列$\{a_n\}$满足$a_1 = 12$，$a_1 + 2a_2 + 3a_3 + \cdots + na_n = n^2a_n$，则$a_{2025} = $\blkda{$\dfrac{4}{675}$}
        \begin{jieda}
            $\left \{ \begin{array}{l}
                 a_1 + 2a_2 + 3a_3 + \cdots + (n-1)a_{n-1} = (n-1)^2a_{n-1}  \\
                 a_1 + 2a_2 + 3a_3 + \cdots + (n-1)a_{n-1} = (n^2-n)a_n
            \end{array} \right.$，因此$\dfrac{a_n}{a_{n-1}} = \dfrac{n-1}{n}$，故累乘得到$a_n = \dfrac{12}{n} (n \geq 2)$，经检验$a_n = \dfrac{12}{n}$，因此$a_{2025} = \dfrac{12}{2025} = \dfrac{4}{675}$
        \end{jieda}
        \item 已知数列$\{a_n\}$的各项均为正数，且$a_1 = 1, (n+1)a_{n+1}^2 - na_n^2 + a_{n+1}a_n = 0$，则$a_n = $\blkda{$\dfrac{1}{n}$}
        \begin{jieda}
            $(n+1)a_{n+1}^2 - na_n^2 + a_{n+1}a_n = 0 \Rightarrow (n+1)\left(\dfrac{a_{n+1}}{a_n}\right)^2 + \dfrac{a_{n+1}}{a_n} - n = 0 \Rightarrow \dfrac{a_{n+1}}{a_n} = \dfrac{n}{n+1} \mbox{或} -1$，因为数列$\{a_n\}$的各项均为正数，所以$\dfrac{a_{n+1}}{a_n} = \dfrac{n}{n+1}$，累乘得$a_n = \dfrac{1}{n} (n \geq 2)$，经检验$a_n = \dfrac{1}{n}$.
        \end{jieda}
        \item 若数列$\{a_n\}$满足$a_n > 0, a_1 = 1, \dfrac{a_{n+1}^2 + a_n^2}{a_{n+1}^2 - a_n^2} = 2n$，则$a_{1013} = $\blkda{45}
        \begin{jieda}
            $\dfrac{a_{n+1}^2 + a_n^2}{a_{n+1}^2 - a_n^2} = 2n$，即$a_{n+1}^2 + a_n^2 = 2n(a_{n+1}^2 - a_n^2)$，即$(2n+1)a_n^2 = (2n-1)a_{n+1}^2$，因此$\left ( \dfrac{a_{n+1}}{a_n}\right)^2 = \dfrac{2n+1}{2n-1}$，因为$a_n > 0$，故$\dfrac{a_{n+1}}{a_n} = \sqrt{\dfrac{2n+1}{2n-1}}$，故累乘得$a_n = \sqrt{2n-1} (n \geq 2)$，经检验$a_n = \sqrt{2n-1}$，故$a_{1013} = \sqrt{2025} = 45$.
        \end{jieda}
        % \item 已知数列$\{a_n\}$满足$a_1 = 1$，$a_{n+1} = \dfrac{a_n}{1+\sqrt{a_n}}(n \in \mathbf{N^*})$，则其前$n$项和$S_n$满足\dotfill{(\kh{A})}
        % \xx{$\dfrac{3}{2} < S_{100} < 3$}{$3 < S_{100} < 4$}{$4 < S_{100} < \dfrac{9}{2}$}{$\dfrac{9}{2} < S_{100} < 5$}
        % \begin{jieda}
        %     根据$y = x$和$y = \dfrac{x}{1+\sqrt{x}}$的关系可知$\{a_n\}$为递减数列，且有$a_n \in (0, 1) (n \geq 2)$，由于$a_2 = \dfrac{1}{2}$，故必有$S_{100} > \dfrac{3}{2}$ \\
        %     考虑倒数法有$\dfrac{1}{a_{n+1}} = \dfrac{1}{a_n} + \dfrac{1}{\sqrt{a_n}} = (\dfrac{1}{\sqrt{a_n}}+\dfrac{1}{2})^2 - \dfrac{1}{4} \Rightarrow \dfrac{1}{\sqrt{a_{n+1}}} - \dfrac{1}{\sqrt{a_n}} < \dfrac{1}{2}$，故累加得$\dfrac{1}{\sqrt{a_n}} < \dfrac{n+1}{2}$($n \geq 2$)，故$a_n \geq \left ( \dfrac{2}{n+1} \right)^2$（当且仅当$n = 1$时等号成立）.代入$a_{n+1} = \dfrac{a_n}{1+\sqrt{a_n}}$有$\dfrac{a_{n+1}}{a_n} = \dfrac{1}{1+\sqrt{a_n}} \leq \dfrac{n+1}{n+3}$（当且仅当$n = 1$时等号成立），故累乘得$a_n \leq \dfrac{6}{(n+1)(n+2)} = 6\left [\dfrac{1}{n+1} - \dfrac{1}{n+2} \right]$（当且仅当$n = 1$时等号成立），因此$S_n < 6(1-\dfrac{1}{n+2})$，$S_{100} < 6(\dfrac{1}{2} - \dfrac{1}{n+2}) < 3$
        % \end{jieda}
    \end{enumerate}
    \item $a_n = S_n - S_{n-1} (n \geq 2)$
    \begin{enumerate}
        \item 已知数列$\{a_n\}$的前$n$项和$S_n = 2 + (n-1)\cdot \left ( \dfrac{1}{2} \right )^{n-1}$，则存在数列$\{x_n\}$，$\{y_n\}$，使得\dotfill{(\kh{D})}
        \xx{$a_n = x_n + y_n$，$n \in \mathbf{N^*}$，其中$\{x_n\}$，$\{y_n\}$是等差数列}
        {$a_n = x_n + y_n$，$n \in \mathbf{N^*}$，其中$\{x_n\}$是等差数列，$\{y_n\}$是等比数列}
        {$a_n = x_n \cdot y_n$，$n \in \mathbf{N^*}$，其中$\{x_n\}$，$\{y_n\}$是等差数列}
        {$a_n = x_n \cdot y_n$，$n \in \mathbf{N^*}$，其中$\{x_n\}$是等差数列，$\{y_n\}$是等比数列}
    \begin{jieda}
        $a_1 = S_1 = 2, a_n = S_n - S_{n-1} = n \cdot \left ( \dfrac{1}{2} \right )^{n-2} (n \geq 2)$，综上，$a_n = n \cdot \left ( \dfrac{1}{2} \right )^{n-2}$\\
        故$x_n = n$，$y_n = \left ( \dfrac{1}{2} \right )^{n-2}$
    \end{jieda}
    \item 设$S_n$为数列$\{a_n\}$的前$n$项和，已知$a_2 = 1, 2S_n = na_n$，则$a_n = $\blkda{$n-1$}
        \begin{jieda}
            $\left \{ \begin{array}{l}
                 2S_n = na_n  \\
                 2S_{n-1} = (n-1)a_{n-1} 
            \end{array} \right. \Rightarrow 2a_n = na_n - (n-1)a_{n-1}$，即$(n-2)a_n = (n-1)a_{n-1}$，即$\dfrac{a_n}{n-1} = \dfrac{a_{n-1}}{n-2} = \dfrac{a_2}{1} = 1$，故$a_n = n-1 (n \geq 2)$ \\
            当$n = 1$时，$2S_n = na_n$即$2a_1 = a_1$，故$a_1 = 0$，因此$a_n = n-1$. \\
            \ps 也可采用累乘法得到$a_n = n - 1$，本质相同.
        \end{jieda}
    \item 记$S_n$为数列$\{a_n\}$的前$n$项和，已知$a_1 = 1$，$\left \{\dfrac{S_n}{a_n}\right \}$是公差为$\dfrac{1}{3}$的等差数列，则$a_n = $\blkda{$\dfrac{n(n+1)}{2}$}
        \begin{jieda}
            因为$\dfrac{S_1}{a_1} = 1$，因此$\dfrac{S_n}{a_n} = \dfrac{n+2}{3}$，即$\left \{ \begin{array}{l}
                 S_n = \dfrac{n+2}{3}a_n  \\
                 S_{n-1} = \dfrac{n+1}{3}a_{n-1}
            \end{array} \right . \\
            \Rightarrow a_n = \dfrac{n+2}{3}a_n - \dfrac{n+1}{3}a_{n-1}$，即$(n-1)a_n = (n+1)a_{n-1}, \dfrac{a_n}{a_{n-1}} = \dfrac{n+1}{n-1}$，累乘得$a_n = \dfrac{n(n+1)}{2}a_1 = \dfrac{n(n+1)}{2} (n \geq 2)$，经检验$a_n = \dfrac{n(n+1)}{2}$
        \end{jieda}
    \item 已知数列$\{a_n\}$的前$n$项和为$S_n$，且$S_n = 1 - a_n$，则数列$\{a_n S_n\}$的前$n$项和$T_n = $\blkda[3]{$\dfrac{2}{3} + \dfrac{1}{3\cdot4^n} - \dfrac{1}{2^n}$}
    \begin{jieda}
        首先可以判断出来$a_n = \dfrac{1}{2^n}$(推导如下： $a_n = (1-a_n) - (1-a_{n-1}) = a_{n-1} - a_n$， 故$a_{n-1} = 2a_n, a_1 = 2$，故$a_n = \dfrac{1}{2^n}$) \\
        故$a_n S_n = \dfrac{1}{2^n} - \dfrac{1}{4^n}$，故$T_n = \left[1- \dfrac{1}{2^n}\right] - \dfrac{1}{3}\left[1-\dfrac{1}{4^n}\right] = \dfrac{2}{3} + \dfrac{1}{3\cdot4^n} - \dfrac{1}{2^n}$
    \end{jieda}
    \item 已知数列$\{a_n\}$的前$n$项和为$S_n$，$a_1 = 1, S_{n+2} = 3S_n + a_{n+2} + 2$，则$a_n = $\blkda[3]{$\left \{\begin{array}{ll}
             4 \cdot 3^{n-2} & ,n \geq 2  \\
             1 & ,n = 1
        \end{array} \right .$}
    \begin{jieda}
        $S_{n+2} = 3S_n + a_{n+2} + 2$即$S_{n+1} = 3S_n + 2$，设$b_n = S_n + 1$，故$b_{n+1} = 3b_n$，即$b_n$是公比为3的等比数列，$b_1 = S_1 + 1 = a_1 + 1 = 2$，故$b_n = 2 \cdot 3^{n-1}$，故$S_n = 2 \cdot 3^{n-1} - 1$，因此$a_n = \left \{\begin{array}{ll}
             4 \cdot 3^{n-2} & ,n \geq 2  \\
             1 & ,n = 1
        \end{array} \right .$
    \end{jieda}
    \item 若数列$\{a_n\}$满足$a_1 + \dfrac{a_2}{2} + \dfrac{a_3}{3} + \cdots + \dfrac{a_n}{n} = 3^n - 2$，则$a_n = $\blkda[3.5]{$\left \{\begin{array}{ll}
             2n \cdot3^{n-1} & ,n \geq 2  \\
             1 & ,n = 1
        \end{array} \right .$}
    \begin{jieda}
        设$b_n = \dfrac{a_n}{n}$，首先计算可知$a_1 = 1, b_1 = 1$. 设$b_n$的前$n$项和为$S_n$，则$\left \{ \begin{array}{l}
             S_n = 3^n - 2  \\
             S_{n-1} = 3^{n-1} - 2
        \end{array} \right. \Rightarrow b_n = \left \{\begin{array}{ll}
             2 \cdot 3^{n-1} & ,n \geq 2  \\
             1 & ,n = 1 
        \end{array} \right .$，故$a_n = \left \{\begin{array}{ll}
             2n \cdot3^{n-1} & ,n \geq 2  \\
             1 & ,n = 1 
        \end{array} \right .$
    \end{jieda}
    \end{enumerate}
    \item 错位相减法求和
    \begin{enumerate}
        \item 记$S_n$为数列$\{a_n\}$的前$n$项和，已知$4S_n = 3a_n + 4$，则$a_n = $\blkda[2]{$4\cdot(-3)^{n-1}$}；设$b_n = (-1)^{n-1}na_n$，则数列$\{b_n\}$的前$n$项和$T_n = $\blkda[3]{$(2n-1)3^n + 1$}
        \begin{jieda}
            $\left \{ \begin{array}{l}
                 S_n = \dfrac{3a_n}{4} + 1  \\
                 S_{n-1} = \dfrac{3a_{n-1}}{4} + 1 
            \end{array}\right. \Rightarrow a_n  = \dfrac{3}{4}(a_n - a_{n-1})$，故$a_n = -3a_{n-1}$，$a_1 = S_1 = \dfrac{3a_1}{4} + 1$，故$a_1 = 4, a_n = 4\cdot(-3)^{n-1}$ \\
            $b_n = 4n \cdot 3^{n-1}$，$\left \{ \begin{array}{l}
                 T_n = 4 \times 1 + 8 \times 3 + \cdots + 4(n-1) \times 3^{n-2} + 4n \times 3^{n-1}  \\
                 3T_n = 4 \times 3 + 8 \times 9 + \cdots + 4(n-1) \times 3^{n-1} + 4n \times 3^n
            \end{array}\right. $ \\
            故$-2T_n = 4 + 4 \times (3 + 9 + \cdots + 3^{n-1}) - 4n \times 3^n = 4 + 4 \times \dfrac{3(1-3^{n-1})}{1-3}  - 4n \times 3^n \\
            = 4 + 2 \times (3^n-3) - 4n \times 3^n = (2-4n)3^n - 2$，因此$T_n = (2n-1)3^n + 1$
        \end{jieda}
        \item 已知等差数列$\{a_n\}$和等比数列$\{b_n\}$均单调递增，前$n$项和分别为$S_n$和$T_n$，且满足$a_1 = b_1, a_3 = b_3, S_3 = 15, T_3 = 14$，则$a_n = $\blkda{$3n-1$}；$b_n = $\blkda{$2^n$}；设$c_n = \dfrac{a_n}{b_n}$，则$\{c_n\}$的前$n$项和$R_n = $\blkda{$5 - \dfrac{5+3n}{2^n}$}
        \begin{jieda}
            $S_3 = 15 \Rightarrow a_2 = 5$，又因为$T_3 = 14, a_1 = b_1, a_3 = b_3$，故$b_2 = 4$，因此$(5-d)(5+d) = 16$，解得公差$d = 3$或$-3$（舍）,故$a_n = 3n-1$，$b_1 = a_1 = 2, b_3 = a_3 = 8$，故公比$q = 2$，故$b_n = 2^n$. \\
            $c_n = \dfrac{a_n}{b_n} = (3n-1) \cdot (\dfrac{1}{2})^n$，$\left \{ \begin{array}{l}
                 R_n = 2 \times \dfrac{1}{2} + 5 \times \dfrac{1}{4} + \cdots + (3n-4) \times \dfrac{1}{2^{n-1}} + (3n-1) \times \dfrac{1}{2^n}  \\
                 \dfrac{1}{2}R_n = 2 \times \dfrac{1}{4} + 5 \times \dfrac{1}{8} + \cdots + (3n-4) \times \dfrac{1}{2^n} + (3n-1) \times \dfrac{1}{2^{n+1}}
            \end{array}\right. $ \\
            故$\dfrac{1}{2}R_n = 1 + 3 \times (\dfrac{1}{4} + \dfrac{1}{8} + \cdots + \dfrac{1}{2^n}) - (3n-1)\times \dfrac{1}{2^{n+1}} = 1 + 3 \times(\dfrac{1}{2} - \dfrac{1}{2^n}) - (3n-1)\times \dfrac{1}{2^{n+1}} = \dfrac{5}{2} - \dfrac{6 + (3n-1)}{2^{n+1}} = \dfrac{5}{2} - \dfrac{5+3n}{2^{n+1}}$，故$R_n = 5 - \dfrac{5+3n}{2^n}$
        \end{jieda}
    \end{enumerate}
    \item 裂项相消求和
    \begin{enumerate}
        \item 等差型 \\
        数列$\{a_n\}$是首项为3，公差为2的等差数列，设数列$\{\dfrac{1}{a_n a_{n+1}}\}$的前$n$项和为$T_n$，则使得$T_n < \dfrac{1}{M}$恒成立的最大的正整数$M = $\blkda{6}
        \begin{jieda}
            $a_n = 2n+1$，则$b_n = \dfrac{1}{a_n a_{n+1}} = \dfrac{1}{(2n+1)(2n+3)} = \dfrac{1}{2}\left [\dfrac{1}{2n+1} - \dfrac{1}{2n+3}\right]$，故$T_n = \dfrac{1}{2}\cdot \left [\dfrac{1}{3} - \dfrac{1}{2n+3}\right] < \dfrac{1}{6}$
        \end{jieda}
        \item 无理型 \\
        已知等差数列 $ \left\{  {a}_{n}\right\} $ 的前 $ n $ 项和为 $ {S}_{n},{a}_{3} + {a}_{5} = {18},{S}_{6} = {48} $ .
        \begin{enumerate}
            \item 求 $ \left\{  {a}_{n}\right\} $ 的通项公式
            \item 设 $ {b}_{n} = \dfrac{2}{\sqrt{{a}_{n + 1}} + \sqrt{{a}_{n - 1}}} $ ,数列 $ \left\{  {b}_{n}\right\} $ 的前 $ n $ 项和为 $ {T}_{n} $ ,证明: 当 $ n \geq  3,n \in  \mathbf{Z} $ 时, $ 4{T}_{n}^{2} > {a}_{n} $ .
        \end{enumerate}
        \begin{jieda}
            \begin{enumerate}
                \item 由题可知, $ \left\{  \begin{array}{l} 2{a}_{1} + {6d} = {18} \\  6{a}_{1} + {15d} = {48} \end{array}\right. $ ,解得 $ \left\{  {\begin{array}{l} {a}_{1} = 3 \\  d = 2 \end{array},\therefore {a}_{n} = {2n} + 1}\right. $ 
                \item  $ {b}_{n} = \dfrac{2}{\sqrt{{a}_{n + 1}} + \sqrt{{a}_{n - 1}}} = \dfrac{2}{\sqrt{{2n} + 3} + \sqrt{{2n} - 1}} = \dfrac{\sqrt{{2n} + 3} - \sqrt{{2n} - 1}}{2} $ 
                
                ${T}_{n} = \\
                \dfrac{1}{2}\left\lbrack  {\left( {\sqrt{5} - 1}\right)  + \left( {\sqrt{7} - \sqrt{3}}\right)  + \left( {\sqrt{9} - \sqrt{5}}\right)  + \cdots  + \left( {\sqrt{{2n} + 1} - \sqrt{{2n} - 3}}\right)  + \left( {\sqrt{{2n} + 3} - \sqrt{{2n} - 1}}\right) }\right\rbrack \\
                = \dfrac{1}{2}\left\lbrack  {\sqrt{{2n} + 1} + \sqrt{{2n} + 3} - 1 - \sqrt{3}}\right\rbrack  , $
                
                当$n \geq  3,n \in  \mathbf{Z}$时，$2{T}_{n} - \sqrt{{a}_{n}} = \sqrt{{2n} + 3} - 1 - \sqrt{3} \geq  2 - \sqrt{3} > 0 $，因此$4{T}_{n}^{2} > {a}_{n} $
            \end{enumerate}
        \end{jieda}
        \item 指数型 \\
        若$a_n = \dfrac{n\cdot2^n}{(n+1)(n+2)}$，则数列$\{a_n\}$的前$n$项和$S_n = $\blkda[2]{$\dfrac{2^{n+1}}{n+2} - 1$}
        \begin{jieda}
            $a_n = \dfrac{2^{n+1}}{n+2} - \dfrac{2^n}{n+1}$，则$S_n = \dfrac{2^{n+1}}{n+2} - 1$
        \end{jieda}
        \item 通项裂项为“+”型 \\
        若$a_n = (-1)^n \cdot \dfrac{2n+1}{n(n+1)}$，则数列$\{a_n\}$的前$n$项和$S_n = $\blkda[4]{$\left \{ \begin{array}{ll}
                 - \dfrac{n+2}{n+1} & n\mbox{为奇数} \\
                 - \dfrac{n}{n+1} & n\mbox{为偶数}
            \end{array}\right. $}
        \begin{jieda}
            $a_n = (-1)^n \cdot \left [\dfrac{1}{n} + \dfrac{1}{n+1}\right]$，因此$S_n = -\left(\dfrac{1}{1} + \dfrac{1}{2}\right) + \left(\dfrac{1}{2} + \dfrac{1}{3}\right) - \cdots + (-1)^n \cdot \left [\dfrac{1}{n} + \dfrac{1}{n+1}\right]$，故当$S_n = \left \{ \begin{array}{ll}
                 -1 - \dfrac{1}{n+1} & n\mbox{为奇数} \\
                 -1 + \dfrac{1}{n+1} & n\mbox{为偶数}
            \end{array}\right. = 
            \left \{ \begin{array}{ll}
                 - \dfrac{n+2}{n+1} & n\mbox{为奇数} \\
                 - \dfrac{n}{n+1} & n\mbox{为偶数}
            \end{array}\right. $ 
        \end{jieda}
    \end{enumerate}
\end{enumerate}